Data assimilation (DA) estimates unknown states or parameters in dynamical systems by combining prior model predictions with noisy observations, typically through variational or Bayesian methods that approximate posterior distributions from finite samples \cite{evensen_data_2022}. Within Bayesian DA, ensemble Kalman filters are efficient but assume Gaussianity and handle only weakly nonlinear models, while particle filters impose few restrictions but need many samples. This trade-off motivates alternatives such as neural-network surrogates for expensive forward models \cite{mucke_deep_2024, cheng_generalised_2022, zhong_reduced-order_2023, chen_reduced-order_2023}, though these deterministic approximations require a priori model-error distributions, hard to specify for high-dimensional stochastic systems.

Generative models offer an alternative, learning to sample directly from a distribution of interest; in DA they are trained on the prior and modified at sampling time to incorporate observations. Two flow-based families dominate: stochastic interpolants (SIs) \cite{albergo_stochastic_2023, albergo_building_2023, chen_probabilistic_2024}, which transport samples through a stochastic differential equation (SDE), and flow matching (FM) \cite{lipman_flow_2023}, which does so through a deterministic ordinary differential equation (ODE); denoising diffusion \cite{ho_denoising_2020, song_score-based_2021} corresponds to SDE sampling of an FM-type model. All three are increasingly applied to physical forecasting, yet methods for turning them into \emph{posterior} samplers have been developed model by model, each tied to whether the generator is stochastic or deterministic.

\paragraph{A unified view.}
We argue that, for data assimilation, this distinction is a choice of sampler within one method, not a difference in method. Stochastic interpolants, flow matching, and denoising diffusion are all special cases of a single interpolant path \cite{albergo_stochastic_2023, lipman_flow_2023, ma_sit_2024, holderrieth_generator_2024}; this unification of the \emph{priors} is established, and our contribution is to use it as a single interface for \emph{posterior sampling}. Conditioning the path on the observation tilts its target to the posterior while leaving the path structure intact, so the conditional score and velocity differ from their prior counterparts by the same likelihood-score term \cite{albergo_stochastic_2023, ma_sit_2024}. Since the conditioned path again admits a one-parameter family of equal-marginal dynamics indexed by a diffusion coefficient \cite{song_score-based_2021}, a single guidance term yields a whole family of provably exact posterior samplers, with three members of practical interest: the SI-SDE (native diffusion), the DM-SDE (a denoising-diffusion sampler lifted from the FM prior), and a deterministic guided ODE \cite{yan_fig_2024, pokle_training-free_2024}. All three share one closed-form observation-interpolant likelihood and differ only in the diffusion coefficient (Figure~\ref{fig:visual_abstract}). Treating them as one family lets us state which efficiency/accuracy trade-offs matter when -- and when the samplers coincide so the cheapest choice suffices.

\input{figures/visual_abstract}

\paragraph{Contributions.}
We introduce a posterior sampling method for data assimilation that applies to any flow-based generative model defined by an interpolant path, without retraining:
\begin{itemize}
    \item \textbf{Unification.} A model-agnostic construction that conditions the interpolant path on the observation and yields, from a single guidance term with weight $\gweight_\tau=\vscoef_\tau+\tfrac12\gdiff_\tau^2$, three posterior samplers -- SI-SDE, DM-SDE, and a deterministic guided ODE -- as the native, lifted, and zero-diffusion members of one family.
    \item \textbf{Theoretical soundness.} We prove the conditioned path is again an interpolant path whose score and velocity correct the prior ones by one likelihood-score term, and that every member of the family generates exact posterior samples in the ideal setting (Theorem~\ref{theorem:unified_posterior}); we also identify the guidance weight consistent with our likelihood surrogate via Doob's $h$-transform.
    \item \textbf{Closed-form likelihood.} An interpolated-observation likelihood whose moments follow from the trained model by one shared Tweedie-type formula, giving a likelihood score with a $\Pi$GDM-style inflated covariance \cite{song_pseudoinverse_2023, chung_diffusion_2023}; this inflation drives accuracy on sparse observations.
    \item \textbf{Computational feasibility.} A Jacobian-free covariance for dense observations, and an ensemble-shared Jacobian approximation reducing the inflated covariance from $\mathcal{O}(E\,N_y)$ network Jacobians per step to a single one, making the accurate sparse-observation regime feasible at field scale.
\end{itemize}
